\chapter{Conclusão}

O presente trabalho abordou o desenvolvimento de um protótipo de controlador eletrônico de velocidade (ESC) para motores de corrente contínua sem escovas (BLDC) utilizando o microcontrolador ESP32. O objetivo geral de estudar o princípio de funcionamento desses conversores e desenvolver uma plataforma física e lógica capaz de oferecer controle de velocidade e torque foi \textbf{parcialmente atingido}. Constatou-se que a arquitetura conceitual, o dimensionamento eletrônico preliminar e a implementação do \textit{firmware} alcançaram maturidade de engenharia consistente e foram validados em bancada. Contudo, a validação dinâmica integral com o motor em rotação não pôde ser concluída em razão de limitações físicas estruturais diagnosticadas no estágio de potência.

A metodologia adotada comprovou a viabilidade da topologia em malha aberta operando com comutação trapezoidal de seis passos. A arquitetura de \textit{firmware} modular demonstrou robustez, isolando as rotinas críticas de controle das tarefas secundárias. Os ensaios estáticos confirmaram a correta manipulação do periférico MCPWM, a eficácia do tempo morto (\textit{dead-time}) implementado por \textit{hardware} e o tempo de resposta adequado das proteções de sobrecorrente (OCP) e subtensão (UVLO).

O limite do escopo prático do protótipo foi estabelecido durante o ensaio a vazio, no qual a ponte inversora apresentou disparo da proteção da fonte primária devido ao consumo transiente excessivo de corrente sem carga acoplada. O diagnóstico metódico deste sintoma isolou a falha no estágio de potência físico, indicando a presença de um caminho de baixa impedância ou degradação de semicondutores em decorrência do estresse de chaveamento, exonerando a camada lógica de controle. Essa capacidade de rastreamento de falhas, evidenciada pela separação clara entre as limitações do \textit{hardware} prototipado e a estabilidade do \textit{software} desenvolvido, atesta a validade do processo investigativo adotado.

Para iterações futuras, a robustez do \textit{hardware} deve ser elevada mediante a adição de isolamento galvânico na interface lógica e a implementação de redes \textit{Snubber} com diodos TVS para a supressão de transientes no barramento. No que tange à integração, a arquitetura de acionamento requer a transição de \textit{drivers} discretos para \textit{Smart Gate Drivers} trifásicos, assegurando proteção nativa e absoluta contra intercondução (\textit{shoot-through}).

No aspecto da manufatura, a prova de conceito em tecnologia de furo passante (THT) cumpriu plenamente seu propósito de instrumentação e depuração inicial. O próximo estágio exige a migração para uma placa de circuito impresso (PCB) multicamadas estritamente SMD (\textit{Surface Mount Device}). Essa transição é essencial para minimizar as indutâncias parasitas, reduzir as sobretensões severas e atenuar a emissão de ruído eletromagnético inerente ao chaveamento em 20 kHz.

Com o \textit{hardware} de potência estabilizado, o \textit{firmware} já dispõe da infraestrutura base, incluindo o módulo de detecção de cruzamento por zero (ZCD) da força contraeletromotriz já codificado, para a validação da malha fechada \textit{sensorless}. Como passo subsequente, sugere-se a evolução do sistema para o Controle Vetorial Orientado a Campo (FOC) e a adoção do barramento CAN (\textit{Controller Area Network}) em substituição à telemetria Wi-Fi, adequando o ESC aos padrões de confiabilidade exigidos pela indústria automotiva e aeroespacial.

Em síntese, o trabalho entrega uma base arquitetural sólida, modular e extensamente documentada. As limitações práticas foram identificadas de forma analítica, consolidando um projeto que não se encerra em uma falha de bancada, mas estabelece um roteiro para a construção futura de controladores BLDC embarcados de alto desempenho.